\section{问题四参数、复现数据与程序}\label{app:q4repro}
\subsection{参数与连续证书细分}
\begin{table}[!htbp]\centering
\caption{问题四最终完整覆盖策略的主要参数}
\begin{tabular}{@{}p{.34\textwidth}p{.59\textwidth}@{}}\toprule
参数 & 取值与作用\\\midrule
连续覆盖布局 & 原点、998 m 八点环、1867 m 十二点环（相位 $15^\circ$）\\
规划用覆盖网格 & 40 m 方格、48 个带角余量的方向区间\\
最终无源核验 & 连续圆弧精确合并；初始 40 m，仅细分未证明格，最细允许 0.625 m\\
角楔与圆外包 & 半宽 $(1.005+10^{-8})^\circ$，每圆 32 条外切半平面\\
定位横向基线 & 前进比例 0.6、侧移上限 100 m、下限 22 m；缩放 $1,1/2,1/4$\\
逼近及清除覆盖 & 最多 8 轮；小区域阈值 75 m；清除覆盖方格间距 28 m\\
单方位沿站补测 & 与上次正测站相距至少 80 m、距估计中心不超过 1450 m\\
后验中心 & 距离 150 点、接收半径 9 点、相容朝向 36 点，全向先验 0.5\\
航路局部搜索 & 上轮热启动加至多 5 个最近邻起点，至多 4 轮翻转与移位\\
信息参考状态 & 1024 个；独立种子 170120；代价下限系数 0.35\\
虚拟时钟上限 & 360000 s；预算不足不宣称完成认证\\\bottomrule
\end{tabular}\end{table}

定位圆外包和双侧负观测截断保持真源包含性，后验中心与信息代价只影响动作选择。实际记录的发射方向字段采用“测点指向源”的比较约定；\cref{eq:q4-reception} 使用物理发射方向，作图时将该字段转过 $180^\circ$。覆盖证明遍历全部连续朝向，因而不依赖字段正反约定。

\begin{table}[!htbp]\centering
\caption{21 站布局的连续证书逐层核验}
\begin{tabular}{@{}rrr@{}}\toprule
方格边长/m & 本层待证明方格 & 核验后未证明方格\\\midrule
40 & 6536 & 348\\
20 & 1072 & 380\\
10 & 1112 & 252\\
5 & 624 & 172\\
2.5 & 336 & 0\\\bottomrule
\end{tabular}\end{table}
后续每层只处理未证明格的子格，并删除与靶区不相交的子格。接受条件一直是整格与全部连续朝向的覆盖。某层未通过只表示未证明，不能仅凭粗网格将其判为真实盲区；源码中的点检仅用于提前拒绝明显不可行布局，不用于接受布局。

\subsection{固定实验定义与统计口径}
四类场景按随机混合、最小半径、贴边朝外、全向编号 $s=0,1,2,3$，第 $i$ 组种子为 $2026091201+1000000s+79i$，$i$ 从零开始。随机混合场景 30 组，其余各 10 组。源数在 $10\sim16$ 离散均匀抽取；随机混合场景中每源以概率 0.5 为定向，方向均匀。最小半径组全部定向且接收半径固定 1000 m；贴边组半径均匀取 $1750\sim1800$ m，接收半径 1000 m，发射半平面朝径向外侧；全向组作为对照。一般源位按圆域面积均匀、接收半径均匀取 $1000\sim1500$ m。同一源、同一测点误差固定，三策略的场景参数完全一致。

后验中心与信息代价各作一项单因素消融，快速策略取 \textrm{fast99}、\textrm{fast95}、\textrm{fast90}，均预先固定随机混合场景前 10 组，复用对应完整覆盖策略作为配对基线。所有 230 次运行均纳入统计；图中代表轨迹固定为随机混合场景第一组，不按表现选例。场景和算法源码校验值保存在 \path{Q4/results/manifest.json}。

bootstrap 重抽样 10000 次，以整组场景为抽样单位，保留策略间配对关系。版本节省率为 $1-\overline T_{\rm new}/\overline T_{\rm old}$；消融增长率为 $\overline T_{\rm ablated}/\overline T_{\rm complete}-1$，区间均为百分位区间，未作多重比较校正。每源时间先算各组比值再平均，不等同于总时间除以总清除数。退化的零失败样本区间不构成未来全清概率为 100\% 的证据。

\begin{table}[!htbp]\centering
\caption{单因素消融：相同十组场景中的完整时间变化}
\begin{tabular}{@{}lrrrr@{}}\toprule
设置 & 时间/s & 增长/\% & 95\% 区间/\% & 全清\\\midrule
关闭后验中心 & 6131.7 & +7.97 & [2.51,13.41] & 10/10 \\
关闭信息代价 & 5826.0 & +2.59 & [-1.47,7.15] & 10/10
\\
\bottomrule\end{tabular}\end{table}

\begin{table}[!htbp]\centering
\caption{快速策略：相同十组场景中的清除结果与两种时间口径}
\begin{tabular}{@{}lrrrrr@{}}\toprule
设置 & 已清/源数 & 全清 & 完整时间/s & 末清每源/s & 退出每源/s\\\midrule
compact & 133/133 & 10/10 & 5679.0 & 392.2 & 436.2 \\
fast99 & 133/133 & 10/10 & 5407.1 & 395.2 & 413.2 \\
fast95 & 130/133 & 8/10 & 4789.7 & 353.1 & 372.0 \\
fast90 & 129/133 & 6/10 & 4210.5 & 323.4 & 330.0
\\
\bottomrule\end{tabular}\end{table}

\subsection{完整场景统计与计算开销}
下表“末清每源”截止末次成功清除，“退出每源”进一步计入残留验证；计算时间为本地进程中策略执行与终止认证的墙钟耗时，不含构造初始化。
\begingroup\small
\begin{longtable}{@{}lrrrrrrr@{}}
\caption{问题四三代策略在全部场景上的补充统计}\\
\toprule 场景/策略 & 已清/源数 & 全清 & 时间/s & 末清每源 & 退出每源 & 检测 & 计算/s\\\midrule\endfirsthead
\toprule 场景/策略 & 已清/源数 & 全清 & 时间/s & 末清每源 & 退出每源 & 检测 & 计算/s\\\midrule\endhead
随机混合/grid & 388/388 & 30/30 & 8426.1 & 552.5 & 666.7 & 237.2 & 0.09 \\
随机混合/joint & 388/388 & 30/30 & 6556.9 & 479.9 & 520.1 & 275.5 & 0.63 \\
随机混合/compact & 388/388 & 30/30 & 5819.1 & 423.7 & 461.3 & 223.0 & 4.19 \\
最小接收半径/grid & 126/126 & 10/10 & 8604.9 & 639.8 & 703.9 & 284.1 & 0.09 \\
最小接收半径/joint & 126/126 & 10/10 & 7535.2 & 602.0 & 614.8 & 329.7 & 0.52 \\
最小接收半径/compact & 126/126 & 10/10 & 6612.2 & 509.4 & 540.1 & 264.7 & 2.43 \\
贴边朝外/grid & 128/135 & 4/10 & 13185.4 & 945.6 & 1032.6 & 449.9 & 0.09 \\
贴边朝外/joint & 135/135 & 10/10 & 8225.0 & 599.2 & 624.9 & 391.6 & 0.44 \\
贴边朝外/compact & 135/135 & 10/10 & 7603.1 & 571.2 & 574.7 & 381.4 & 1.07 \\
全向/grid & 130/130 & 10/10 & 7437.9 & 373.1 & 580.7 & 197.7 & 0.06 \\
全向/joint & 130/130 & 10/10 & 6013.8 & 438.6 & 473.7 & 257.9 & 0.50 \\
全向/compact & 130/130 & 10/10 & 5328.0 & 383.0 & 419.9 & 205.5 & 3.24
\\
\bottomrule\end{longtable}\endgroup

原始离线实验采用 Python 3.12.14、NumPy 2.3.5，六个独立进程运行，单局策略执行墙钟最大 7.55 s。并行资源争用和本机环境会改变墙钟时间，不能将其替代正式模拟器的在线程序运行时间。所有 230 份完整轨迹、65229 条动作通过独立落盘复核，移动、检测、切频、成功及失败清除的账本均吻合，拒绝动作数为零。六组已有专项检查共 37 项通过，包括连续证书、极端定位、航路及整策略回归；日志见 \path{Q4/results/}。

本次附件整理另在 macOS 上进行了 17 次试运行，两个要求连续认证的策略共 8 次均全清并认证；其中部分策略的路径及虚拟时间与历史 Windows 记录不同，详见 \path{Q4/results/package_validation.json}。因此，保存动作的统计可独立重算，而跨平台重新运行不保证逐动作一致；本文统计仍取原始固定实验，不混入新结果。

\subsection{支撑材料索引}
附件目录、数据来源及四类运行命令统一见附录~\ref{app:files}。
逐组结果与完整动作记录分别存于 \path{Q4/results/} 和 \path{Q4/results/traces/}；
正文图、完整补充图册分别为 \path{Q4/figures/} 和 \path{Q4/supplement.pdf}。
运行环境与已完成核验见 \path{Q4/README.md}，论文数字及图表的证据对应见
\path{Q4/DATA_SOURCES.md}。各程序均通过 \path{Q4/reproduce.py} 调用。

\section{问题四最终策略核心代码}\label{app:q4code}
以下直接载入最终策略的核心源码，展示紧凑站位、联合航路与后验估价的接入方式，以及根据实际负观测逐频道核验完成证书的逻辑。
自适应认证、航路优化、后验估计和定位清除分别由
\path{Q4/code/p4_layout_v2.py}、\path{p4_route_v2.py}、\path{p4_posterior.py}
和 \path{p4_homing.py} 实现；三代策略、全部依赖及运行入口均保留在 \path{Q4/} 附件中。
\lstinputlisting[language=python, caption={Q4/code/p4\_compact.py：最终紧凑策略}]{../attachments/Q4/code/p4_compact.py}
